% Section 8: The dual} | Part: 1 | Depends on: see docs/REWRITE_PLAN.md | Status: todo | Original lines: 1555-1596
%!TEX root = ../../../main.tex
\subsection{The dual}

\begin{definition}[Dual and undual]\label{def:dual}
The \emph{dual} of a multivector $A$ is
\[
 \dual A:=A\,\pss^{-1}\;=\;A\lc\pss^{-1},
\]
and the \emph{undual} is $A\mapsto A\pss$. (The two operations differ from each other only by the scalar factor $\pss^{2}$: $A\pss=\pss^2\,\dual A$.)
\end{definition}

\begin{theorem}[Properties of the dual]\label{thm:dual}
\begin{enumerate}
  \item $A\mapsto\dual A$ is a linear isomorphism $\GA^k(V)\to\GA^{n-k}(V)$, and $\dual{(\dual A)}=\pss^{-2}A=\pm A$.
  \item (\emph{Orthogonal complement}) If $A$ is a $k$-blade representing the subspace $U\subset V$, then $\dual A$ is an $(n-k)$-blade representing the orthogonal complement $U^\perp$ and $\norm{\dual A}=\norm A$ (Euclidean case).
  \item (\emph{Wedge $\leftrightarrow$ contraction}) For all multivectors $A,B$,
  \begin{equation}\label{eq:dualduality}
   \dual{(A\wedge B)}=A\lc\dual B,\qquad \dual{(A\lc B)}=A\wedge\dual B .
  \end{equation}
\end{enumerate}
\end{theorem}
\begin{proof}
(1) follows from Proposition~\ref{prop:pss}(3) and invertibility of $\pss$.
(2) Choose an orthonormal basis adapted to $U$ so that (after scaling) $A=e_1\cdots e_k$. Then $\pss^{-1}=e_n^{-1}\cdots e_{k+1}^{-1}\,e_k^{-1}\cdots e_1^{-1}$ and, since the $(n-k)$-blade $e_n^{-1}\cdots e_{k+1}^{-1}$ is built from vectors orthogonal to $e_1,\dots,e_k$, it commutes with $A$ up to $(-1)^{k(n-k)}$. Hence $A\pss^{-1}=(-1)^{k(n-k)}e_n^{-1}\cdots e_{k+1}^{-1}$, a blade spanning $U^\perp$.
(3) For homogeneous $B$ we have $B\lc\pss^{-1}=B\pss^{-1}$ (the product has the single grade $n-\deg B$). Proposition~\ref{prop:contr}(3) gives
$A\lc(B\pss^{-1})=A\lc(B\lc\pss^{-1})=(A\wedge B)\lc\pss^{-1}=\dual{(A\wedge B)}$, the first identity. For the second, apply the first to $A$ and $B'=B\pss$ (so $\dual{B'}=B$): $A\lc B=\dual{(A\wedge B\pss)}$; then $\dual{(A\lc B)}=(A\wedge B\pss)\pss^{-2}=A\wedge(B\pss^{-1})=A\wedge\dual B$, using that $\pss^{2}$ is a scalar.
\end{proof}

\begin{corollary}[Meet and the regressive product]\label{cor:meet}
Define the \emph{meet} (regressive product)
\[
 A\meet B:=(\dual B\wedge\dual A)\,\pss=\dual B\lc A .
\]
If $A,B$ are blades whose subspaces $U_A,U_B$ have $U_A+U_B=V$, then $A\meet B$ is (a multiple of) the blade of $U_A\cap U_B$.
\end{corollary}
\begin{proof}
$U_A^\perp$ and $U_B^\perp$ span $(U_A\cap U_B)^\perp$ when $U_A+U_B=V$, so $\dual B\wedge\dual A$ represents $(U_A\cap U_B)^\perp$ by Theorem~\ref{thm:dual}(2), and undualizing gives $U_A\cap U_B$. The equality of the two expressions follows from \eqref{eq:dualduality}: $(\dual B\wedge\dual A)\pss=\pss^2\,\dual{(\dual B\wedge\dual A)}=\pss^2\,\dual B\lc\dual{(\dual A)}=\dual B\lc A$.
\end{proof}

\begin{example}[Meet of two planes in $\R^3$]
$A=e_1e_2$ and $B=e_2e_3$ (the $xy$- and $yz$-planes). With $\pss=e_1e_2e_3$, $\pss^{-1}=-\pss$ one finds $\dual A=e_3$, $\dual B=e_1$, and $A\meet B=\dual B\lc A=e_1\lc(e_1e_2)=e_2$: the $y$-axis, the line of intersection.
\end{example}
